\section{Introduction}
% P1 1. metion video generation 2. it is costly
Video diffusion transformers have rapidly advanced in generation fidelity, temporal coherence, and controllability, fueling the growing popularity and adoption of AI-generated video across content creation, entertainment, and online media~\citep{team2025wan,kong2024hunyuanvideo,yang2025cogvideox}.
However, the inference cost of video generation remains substantial, driven by iterative denoising over long spatiotemporal token sequences with large Transformer backbones.
For example, generating a five-second 1080p video using ByteDance's Seedance 2.0 costs approximately \$1.87 through its commercial API~\citep{byteplus2026seedance2}.
Efficient inference has therefore become a key requirement for scalable video generation services.

% P2 previous intra-request attempts
To reduce inference cost, prior work has explored a broad range of intra-request acceleration techniques.
These approaches exploit computational redundancy, including few-step distillation, feature caching~\citep{nie2026transition, lyu2025fastercache}, token-wise computation reduction, and sparse attention~\citep{zou2024accelerating, chen2026sparse}.
Collectively, these methods target the computational cost of individual generation trajectories.
At the service level, video generation systems operate under continuous workloads consisting of requests from many different users.
These requests often exhibit recurring semantic and structural similarities, including objects, scenes, and backgrounds.
This workload pattern reveals a new opportunity for inference acceleration by exploiting computational redundancy across different generation requests. 

% P3 inter-request
Recent studies have begun exploring this inter-request acceleration opportunity through cache reuse across related generation requests.
Existing T2V approaches range from retrieving similar requests and reusing intermediate noise or latent states~\citep{cherian2026chai}, to entity- or region-selective recomputation~\citep{liu2026beyond}.
Going beyond textual requests, recent work on I2V exploits similar attention patterns across globally similar scenes to reuse historical sparse-attention connectivity at inference time~\citep{liu2026chorus}.
% Collectively, these studies establish the potential of inter-request reuse, while underscoring the need for principled compatibility assessment to determine both the validity and extent of historical computation reuse under varying degrees of request similarity.
Collectively, these studies reuse latent and attention caches from previous requests, establishing the potential of inter-request reuse while underscoring the need for principled assessment of reuse validity and extent across requests with varying similarity.

% P4 the gap [TODO] 添加一下青山雪山橄榄球杠铃的对比图
However, existing cross-request methods largely rely on textual semantic similarity to identify reusable histories, which provides only a coarse proxy for generation-level compatibility.
% Even highly similar prompts may produce different backgrounds, colors, or object geometries, while objects with comparable size and position can still exhibit incompatible shapes, making direct reuse unsafe.
For example, replacing an apple rolling down a snowy mountain with a lychee may preserve substantial textual similarity while inducing a different background or color composition, while objects with similar size and position, such as a football and a barbell, may still exhibit incompatible geometry that prevents direct reuse.
Moreover, as cross-request reuse spans both latent trajectories and attention connectivity, reliable acceleration requires increasingly fine-grained compatibility evidence.
Such evidence is often obtained either from richer conditioning inputs such as images or through auxiliary post hoc processing, where external language models and segmentation tools construct region masks for reuse, introducing additional inference overhead.
These observations expose two key gaps in existing cross-request acceleration methods, namely reliable generation-level compatibility assessment and principled determination of reuse extent across latent trajectories and attention connectivity.
% Meanwhile, the evolving internal states of video diffusion models expose object-level structural evidence that remains underexplored for cross-request compatibility assessment.
% [TODO] insight 加到下一段，引用 attn 那篇文章

% P5 Carnator
To bridge these gaps, we present \emph{Carnator}, an inter-request acceleration framework for T2V generation that identifies compatible historical computation and jointly exploits it across latent trajectories and attention connectivity.
\emph{Carnator} first performs a lightweight early probe to construct a generation-native object signature that captures object masks and reuse-relevant visual characteristics directly from the model’s internal generation states, without auxiliary models or post hoc mask generation.
The resulting signature serves as the basis for generation-level compatibility filtering, providing a more reliable criterion for historical computation reuse than textual similarity alone.
Once a cache is accepted, \emph{Carnator} jointly exploits historical computation across both latent trajectories and attention connectivity.
This design improves computational efficiency while preserving the target-specific updates necessary for generation fidelity.

Our main contributions are summarized as follows:
\begin{itemize}[leftmargin=*]
    \item We introduce a generation-native object-aware retrieval mechanism for cross-request cache reuse, which constructs early signatures directly from the internal states of video diffusion models. 
    \emph{Carnator} combines attention-derived object regions with spatial, size, position, and appearance characteristics to identify historical generation states that are more compatible with the target request than prompt similarity alone.

    \item We propose a joint latent-attention reuse mechanism that exploits cross-request redundancy in both latent trajectories and attention connectivity.
    \emph{Carnator} preserves cached latent states outside target-specific active regions while recomputing active regions with contextual information, and reuses historical sparse-connectivity masks in a separate set of Transformer layers while recomputing target attention representations.

    \item We conduct extensive experiments on diverse text-to-video workloads to evaluate both cross-request cache retrieval and inference acceleration.
    \emph{Carnator} improves cache retrieval quality by [X]\%, reduces false reuse by [Y]\%, and achieves up to [Z]$\times$ end-to-end acceleration while maintaining competitive generation quality.
\end{itemize}

